\documentclass[a4paper,11pt]{article}

\usepackage{revy}
\usepackage{amssymb}
\usepackage{amsmath}
\usepackage{commath}
\usepackage{amsfonts}

\revyname{MatematikRevyen}
\revyyear{2025}
\version{0.1}
\eta{? minutter}
\status{Færdig}

\title{Restaurant Dumpet}
\author{Augusta '20, Hugo '21, Sirius '23, Elinborg '22, Theis '22}

\begin{document}
\maketitle

\begin{roles}
\role{I}[KE] Instruktør
\role{S1}[Elinborg] Studerende på SU
\role{S2}[Viktor] Studerende på SU nr. 2
\role{T}[Vitus] Fransk Tjener (evt. Fysiker)
\end{roles}

\begin{props}
\prop{1} HMM
\end{props}


\begin{sketch}


\scene Stearinlys og fin dug på et lille bord. Tjener står afventende - evt. lidt fin klavermusik. To broke SU studerende på date. Et billede af Salvador Dahli i baggrunden.

\says{T} Velkommen! Velkommen til Restauranten Du dumpede et enkelt lille bitte fag på første år og nu er du løbet tør for SU. Ta plads. 

\scene T trækker stolen ud og S1 og S2 sætter sig ned

\says{S} Vi har en dahlig dahlig menu i dag.

\scene Tager en menu frem og peger.

\says{T} Til forret jeg anbefaler vores Soupe Lentille de Tomate - på franske linser naturligvis.

\scene Bordet kigger på hinanden og nikker.

\says{T} Til hovedret - dehydrerede linser i et friskt krydret indkog af tomat og kokosmælk.

\says{S} Er det Røde linser?

\says{T} Naturligvis, naturligvis. Og til dessert monsieurs?

\says{S} I har dessert?

\says{T} Naturligvis, naturligvis. En fin coulis af tomat og kokus - med knas af røde frø fra ærteblomsten ægte linse.

\says{S} Så tre gange dahl?

\says{T} Naturligvis - s'gerne, s'gerne.
\end{sketch}
\end{document}

